\section{Related work}
Our work is largely inspired by the entropy search (ES) methods~\cite{hennig2012,hernandez2014predictive}, which established the information-theoretic view of Bayesian optimization by evaluating the inputs that are most informative to the $\argmax$ of the function we are optimizing. 

Our work is also closely related to probability of improvement (PI)~\cite{kushner1964}, expected improvement (EI)~\cite{mockus1974}, and  the BO algorithms using upper confidence bound to direct the search~\cite{auer2002b,kawaguchi2015bayesian,kawaguchi2016global}, such as  GP-UCB~\cite{srinivas2009gaussian}. In~\cite{wang2016est}, it was pointed out that GP-UCB and PI are closely related by exchanging the parameters. Indeed, all these algorithms build in the heuristic that the next evaluation point needs to be likely to achieve the maximum function value or have high probability of improving the current evaluations, which in turn, may also give more information on the function optima like how ES methods queries. These connections become clear as stated in Section 3.1 of our paper.

Finding these points that may have good values in high dimensional space is, however, very challenging. In the past, high dimensional BO algorithms were developed under various assumptions such as the existence of a lower dimensional function structure~\cite{djolonga2013high,wang2013}, or an additive function structure where each component is only active on a lower manifold of the space~\cite{li2016high,kandasamy2015high}. In this work, we show that our method also works well in high dimensions with the additive assumption made in~\cite{kandasamy2015high}.

